\section*{Capítulo \ref{ch:b3:many-electron-atoms} --- Átomos polieletrônicos e símbolos de termo}

\begin{solution}{exo:b3:many-electron-atoms:1}
\[
\Psi = \frac{1}{\sqrt6}\begin{vmatrix}1s\alpha(1) & 1s\beta(1) & 2s\alpha(1)\\
1s\alpha(2) & 1s\beta(2) & 2s\alpha(2)\\ 1s\alpha(3) & 1s\beta(3) & 2s\alpha(3)
\end{vmatrix}.
\]
$1s$ oferece apenas dois \omterm{def:b3:many-electron-atoms:spin-orbital}{spin-orbitais}, $1s\alpha$ e $1s\beta$; um terceiro elétron
repetiria um deles, tornando duas colunas iguais: o determinante é nulo.
\end{solution}

\begin{solution}{exo:b3:many-electron-atoms:2}
$\binom62 = 15$, $\binom63 = 20$, $\binom{10}{2} = 45$. $p^3$: ${}^4$S (4) +
${}^2$D (10) + ${}^2$P (6) = 20. $d^2$: ${}^3$F (21) + ${}^3$P (9) + ${}^1$G (9) +
${}^1$D (5) + ${}^1$S (1) = 45.
\end{solution}

\begin{solution}{exo:b3:many-electron-atoms:3}
N ($2p^3$, semipreenchida): \termsym{4}{S}{3/2}. O ($2p^4$, mais que semipreenchida):
\termsym{3}{P}{2}. F ($2p^5$): \termsym{2}{P}{3/2}. \ce{Ti^2+} ($3d^2$):
\termsym{3}{F}{2}. \ce{Fe^3+} ($3d^5$, semipreenchida, todos os spins paralelos, $L = 0$):
\termsym{6}{S}{5/2}.
\end{solution}

\begin{solution}{exo:b3:many-electron-atoms:4}
${}^2$D: $J = \frac52$ (6 estados) e $\frac32$ (4); $6 + 4 = 10 = 5 \times 2$.
${}^3$P: $J = 2$ (5), 1 (3), 0 (1); $5 + 3 + 1 = 9 = 3 \times 3$.
\end{solution}

\begin{solution}{exo:b3:many-electron-atoms:5}
Intervalos $16.4$ ($J = 1 - 0$) e $27.0$ ($2 - 1$): razão 1.65 em vez de 2.
$A = 16.4/1 = \qty{16.4}{cm^{-1}}$ pelo primeiro, $27.0/2 = \qty{13.5}{cm^{-1}}$
pelo segundo. Um único $A$ não se ajusta: os níveis estão levemente misturados
com os de outros termos (${}^1$D, ${}^1$S), e o \omterm{def:b3:many-electron-atoms:russell-saunders}{acoplamento de Russell--Saunders}
é uma aproximação.
\end{solution}

\begin{solution}{exo:b3:many-electron-atoms:6}
$\lambda = 10^7/\tilde\nu$: \qty{589.76}{nm} (D$_1$) e \qty{589.16}{nm}
(D$_2$). Separação $\qty{17.20}{cm^{-1}} = \qty{2.13}{meV}$. $E(\frac32) -
E(\frac12) = \frac32A$, logo $A = \qty{11.47}{cm^{-1}}$.
\end{solution}

\begin{solution}{exo:b3:many-electron-atoms:7}
$3p \to 3s$: permitida ($\Delta l = -1$). $3d \to 3s$: proibida ($\Delta l = -2$).
$3d \to 3p$: permitida. $4s \to 3s$: proibida ($\Delta l = 0$, sem mudança de
paridade). $4s \to 3p$: permitida.
\end{solution}

\begin{solution}{exo:b3:many-electron-atoms:8}
Média de ${}^3$P: $(5 \times 169086.76 + 3 \times 169086.84 + 169087.83)/9 =
\qty{169086.91}{cm^{-1}}$. Distância até ${}^1$P: \qty{2047.98}{cm^{-1}} $=
\qty{0.254}{eV}$, logo $K(1s,2p) = \qty{0.127}{eV}$, três vezes menor que
$K(1s,2s)$. A \omterm{def:b3:many-electron-atoms:exchange}{integral de troca} mede a sobreposição das densidades dos dois
orbitais; o orbital $2s$ penetra na região de $1s$ (tem densidade perto do
núcleo), o orbital $2p$ se anula no núcleo e se sobrepõe menos.
\end{solution}

\begin{solution}{exo:b3:many-electron-atoms:9}
O maior $M_L = 4$ exige os dois elétrons em $m_l = 2$, com spins opostos: um
singleto, ${}^1$G (9 estados). O seguinte, $M_L = 3$, vem de $m_l = 2, 1$, com
$M_S = 1$ possível: ${}^3$F (21). Resta $M_L = 2$ apenas com $M_S = 0$:
${}^1$D (5). Resta $M_L = 1$ com $M_S = 1$: ${}^3$P (9). Sobra um estado com
$M_L = M_S = 0$: ${}^1$S. Total 45.
\end{solution}

\begin{solution}{exo:b3:many-electron-atoms:10}
Em $\mathbf r_1 = \mathbf r_2 = \mathbf r$, a função do tripleto é
$\frac{1}{\sqrt2}[1s(\mathbf r)2s(\mathbf r) - 2s(\mathbf r)1s(\mathbf r)] = 0$; a
função do singleto é $\sqrt2\,1s(\mathbf r)2s(\mathbf r) \ne 0$. No tripleto, os
elétrons nunca são encontrados no mesmo ponto e ficam, em média, mais
afastados, de modo que sua repulsão é menor: $E_+ - E_- = 2K > 0$.
\end{solution}

\begin{solution}{exo:b3:many-electron-atoms:11}
$d^8$ está mais que semipreenchida: o acoplamento é invertido, $J = L + S = 4$
é o mais baixo. Intervalos: $E(3) - E(4) = 1360.7$ e $E(2) - E(3) = 908.9$; razão
1.50, contra $4/3 = 1.33$ de Landé. $|A| = 1360.7/4 = \qty{340}{cm^{-1}}$
($A < 0$). O \omterm{def:b3:many-electron-atoms:spin-orbit}{acoplamento spin--órbita} é muito maior nesse íon mais pesado que
no carbono (\qty{16}{cm^{-1}}).
\end{solution}

\begin{solution}{exo:b3:many-electron-atoms:12}
$\sum_J(2J+1)[J(J+1) - L(L+1) - S(S+1)] = 0$ (os $(2L+1)(2S+1)$ estados podem
ser contados em qualquer das duas bases, e $\sum M_LM_S = 0$ sobre os estados desacoplados).
Para ${}^3$P: $1(0 - 4) + 3(2 - 4) + 5(6 - 4) = -4 - 6 + 10 = 0$. A média
ponderada do carbono é $(0 + 3 \times 16.42 + 5 \times 43.41)/9 = \qty{29.6}{cm^{-1}}$:
a energia que o termo teria sem \omterm{def:b3:many-electron-atoms:spin-orbit}{acoplamento spin--órbita}.
\end{solution}

\begin{solution}{pb:b3:many-electron-atoms:1}
\textbf{1.} $1s\alpha$, $1s\beta$, $2s\alpha$, $2s\beta$.
\textbf{2.} $|1s\alpha\,2s\alpha|$, $|1s\alpha\,2s\beta|$, $|1s\beta\,2s\alpha|$,
$|1s\beta\,2s\beta|$.
\textbf{3.} $|1s\alpha\,2s\alpha| = \frac{1}{\sqrt2}[1s(1)2s(2) - 2s(1)1s(2)]\,
\alpha(1)\alpha(2)$, desenvolvendo o determinante $2\times2$.
\textbf{4.} Sua soma é $\frac{1}{\sqrt2}[1s(1)2s(2) - 2s(1)1s(2)]\cdot
\frac{1}{\sqrt2}[\alpha(1)\beta(2) + \beta(1)\alpha(2)]$ (vezes $\sqrt2$, depois
normalizada); sua diferença dá $\frac{1}{\sqrt2}[1s(1)2s(2) +
2s(1)1s(2)]\cdot\frac{1}{\sqrt2}[\alpha(1)\beta(2) - \beta(1)\alpha(2)]$.
\textbf{5.} Parte espacial simétrica com a função de spin antissimétrica ($S = 0$);
parte espacial antissimétrica com as três funções de spin simétricas ($S = 1$).
\textbf{6.} $2S + 1 = 1$ e 3: um e três estados de mesma energia.
\textbf{7.} $[\ce{Ar}]\,3d^2$; $\binom{10}{2} = 45$ determinantes.
\textbf{8.} No máximo $S = 1$ e, no máximo, $L = 2 + 1 = 3$; menos que semipreenchida:
\termsym{3}{F}{2}, o nível em 0 nos dados.
\textbf{9.} $21 + 9 + 9 + 5 + 1 = 45$.
\textbf{10.} Normal ($J = 2$ mais baixo). Intervalos 184.9 e 235.5: razão 1.27
contra $4/3 = 1.33$; a regra vale com 5\,\% de margem.
\textbf{11.} ${}^3$F: $(5 \times 0 + 7 \times 184.9 + 9 \times 420.4)/21 =
\qty{241.8}{cm^{-1}}$; ${}^3$P: $(10538.4 + 3 \times 10603.6 + 5 \times
10721.2)/9 = \qty{10661.7}{cm^{-1}}$.
\textbf{12.} $15B = \qty{10419.9}{cm^{-1}}$, $B = \qty{695}{cm^{-1}}$.
\textbf{13.} $\Delta S = 0$.
\textbf{14.} $\Delta S = 1$ é proibida, e $1s2s \to 1s^2$ tem $\Delta l = 0$
(sem mudança de paridade; ${}^3$S$_1 \to {}^1$S$_0$ também tem $\Delta L = 0$ entre dois estados
S). O estado é metaestável: vive muito mais tempo que os estados excitados
comuns.
\textbf{15.} $169086.91 - 159855.97 = \qty{9230.94}{cm^{-1}}$: \qty{1083.3}{nm}, no
infravermelho próximo.
\textbf{16.} $171134.89 - 166277.44 = \qty{4857.45}{cm^{-1}}$: \qty{2058.7}{nm}.
\textbf{17.} $\qty{171134.89}{cm^{-1}}$: \qty{58.43}{nm}, no ultravioleta de
vácuo (absorvido pelo ar).
\textbf{18.} Cada família tem suas próprias raias e seu próprio estado mais baixo,
e as duas nunca se ligam pela luz: elas se comportam como dois gases com dois espectros.
\textbf{19.} $E_\pm = E_{1s} + E_{2s} + J \pm K$ (singleto $+$, tripleto $-$).
\textbf{20.} O tripleto: sua função espacial se anula quando os elétrons se
encontram, de modo que eles se repelem menos.
\textbf{21.} \qty{6421.47}{cm^{-1}}, $6421.47/8065.54 = \qty{0.796}{eV}$.
\textbf{22.} Distância $\qty{2047.98}{cm^{-1}} = \qty{0.254}{eV}$, logo $K(1s,2p) =
\qty{0.127}{eV}$.
\textbf{23.} O $2s$ penetra na região de $1s$ e se sobrepõe a ela mais do que
o $2p$.
\textbf{24.} Sim: nas duas configurações, o tripleto ($S$ maior) é o mais baixo.
\textbf{25.} $K = \frac12 \times \qty{0.796}{eV}$: \textbf{a \omterm{def:b3:many-electron-atoms:exchange}{integral de troca}
$K(1s,2s)$ do hélio é \qty{0.398}{eV}}, medida como a metade da
separação singleto--tripleto.
\end{solution}
